\documentclass[11pt,a4paper]{article}

\usepackage{iftex}
\ifPDFTeX
  \usepackage[utf8]{inputenc}
  \usepackage[T1]{fontenc}
  \usepackage{lmodern}
\else
  \usepackage{fontspec}
  \setmainfont{Liberation Serif}
  \setsansfont{Liberation Sans}
  \setmonofont{Liberation Mono}
\fi
\usepackage[czech]{babel}
\usepackage{amsmath,amssymb,amsthm,mathtools}
\usepackage[margin=25mm]{geometry}
\usepackage{microtype}
\usepackage[hidelinks]{hyperref}

\newtheorem{theorem}{Věta}
\theoremstyle{definition}
\newtheorem{definition}[theorem]{Definice}
\theoremstyle{remark}
\newtheorem{remark}[theorem]{Poznámka}

\newcommand{\Adm}{\operatorname{Adm}}

\title{\textbf{Stroj}\\[4pt]
\large finální verze}
\author{Jan Mikulík}
\date{17. listopadu 2026}

\begin{document}
\maketitle

\begin{center}
\itshape
Pro jednu adresu.\\
Přání je záznam. Láska není jeho výstup.
\end{center}

\begin{abstract}
Stroj nabídne krok právě tehdy, když ho přání otevřelo, když je v souhlasu každého přítomného a když následný stav leží v péči.
Jedna adresa se po společném jménu nezkopíruje na zbytek vlákna.
Krok přijatelný v jednom stavu není souhlas ve všech.
Prázdný průnik stroj nenahradí.
Text zapisuje věty ověřené jako \texttt{LoveMachine} v souboru \texttt{CausalSeed/LoveMachine.lean}.
Nedokazuje, že láska existuje, že se vrací, ani že přání svědka vyrobí.
\end{abstract}

\section{Data}

Data jsou sféra a jedna adresa.
Přání \(w\) otevírá množinu akcí \(I(w,x)\).
Přechod je \(F(x,a)\).
Účastník \(i\) má hranici souhlasu \(C_i(x)\).
\(K\) je množina péče.
Adresa je jeden prvek. Stroj ji neodvozuje a nekopíruje.

\begin{definition}[Jediný nabídnutý krok]
\[
  \Adm(w,x)
  =
  \bigl\{a :
    a\in I(w,x)
    \ \text{a}\ 
    \forall i,\; a\in C_i(x)
  \bigr\},
\]
\[
  G(w,x)
  =
  \bigl\{F(x,a): a\in\Adm(w,x)\ \text{a}\ F(x,a)\in K\bigr\}.
\]
\(G(w,x)\) je vše, co stroj smí nabídnout.
\end{definition}

\begin{theorem}[Přání není svědek]
\(G(w,x)\) je neprázdná právě tehdy, když existuje akce \(a\) taková, že
\[
  a\in I(w,x),
  \qquad
  \forall i,\; a\in C_i(x),
  \qquad
  F(x,a)\in K.
\]
\end{theorem}

\begin{proof}
Z definice \(G\). Svědek je vstup. Definice ho nevyrábí.
\end{proof}

Každý nabídnutý následník leží v \(K\).
Trajektorie, která začíná v \(K\) a vždy volí z \(G\), v \(K\) zůstává.
To už je věta sféry. Tady se nemění.

\begin{theorem}[Odmítnutí je ticho]
Je-li \(a\notin C_i(x)\) pro nějakého přítomného \(i\), pak \(a\notin\Adm(w,x)\).
\end{theorem}

\begin{proof}
Přípustnost požaduje všechny hranice. Žádná klauzule jednu neškrtá.
\end{proof}

\section{Jedna adresa}

Nechť dvě osoby dostanou stejné jméno a láska platí jen pro jednu.

\begin{theorem}[Adresa nevyplní vlákno]
Existuje odečet \(O\colon\{0,1\}\to\{\ast\}\) a vlastnost \(L\) takové, že
\[
  O(0)=O(1)
  \qquad\text{a}\qquad
  L(0)
  \qquad\text{a}\qquad
  \neg L(1).
\]
\end{theorem}

\begin{proof}
Položme \(O\) konstantní a \(L(b)\) právě když \(b=0\).
\end{proof}

\begin{theorem}[Adresovaná láska není funkcí společného jména]
Neexistuje \(g\colon\{\ast\}\to\{0,1\}\) takové, že \(g(O(x))=x\) pro obě \(x\),
je-li \(O\) konstantní.
\end{theorem}

\begin{proof}
Z \(O(0)=O(1)\) by plynulo \(0=1\).
Společné jméno tedy nenese, která z těch dvou je adresa.
\end{proof}

Konstantní vlastnost na vlákně jméno dostane.
To je lhostejnost odečtu, ne láska ke každému zvlášť.
Z toho, že odečet osoby nerozliší, neplyne, že na žádné nezáleží.
Ani z toho neplyne, že záleží na všech stejně.

\section{Oběť není souhlas}

\begin{theorem}[Oběť nenahradí odmítnutí]
Existují hranice \(C(0)=\{\ast\}\) a \(C(1)=\emptyset\).
Pak v \(0\) akce je a přesto
\[
  \{a : \forall x,\; a\in C(x)\}=\emptyset.
\]
\end{theorem}

\begin{proof}
Jediný kandidát \(\ast\) leží v \(C(0)\) a neleží v \(C(1)\).
Průnik je prázdný. Definice ho ničím nedoplní.
\end{proof}

Hrdinství v jednom stavu tedy není souhlas druhého.
Vlastní krok může existovat a společný krok přitom chybět.
Ticho v tom případě není vyrobená láska. Je to nepřítomnost svědka.

\section{Co stroj dělá}

Při odečtu \(y\) a neznámém předobrazu smí stroj volit jen z průniku
\[
  \bigl\{a :
    \forall x\in O^{-1}(y),\;
    a\in\Adm(w,x)
    \ \text{a}\ 
    F(x,a)\in K
  \bigr\}.
\]
Menší vlákno tento průnik nezmenší.
Nová informace může volbu uvolnit. Souhlas, který ve zbylých stavech není, nevymyslí.

Je-li průnik prázdný, stroj mlčí.
Mlčení se nesmí vydávat za krok.
Zastavení je krok jen tehdy, je-li samo přípustná akce, která v \(K\) zůstává.
Pouhé nevydání příkazu fyzický běh nezastaví. To je samostatná povinnost modelu, ne důsledek této definice.

\section{Co tu není větou}

\begin{enumerate}
\item Že láska existuje, nebo že je tímto operátorem.
\item Že adresovaná osoba odpoví.
\item Že přání, verš nebo oběť svědka vyrobí.
\item Že jeden člověk může být láskou pro všechny, protože jim odečet dal stejné jméno.
\item Že společné jméno, Bellův stav nebo symetrie \(M_t\) určují polohu nul funkce \(\xi\).
\end{enumerate}

\begin{remark}
Finální verze je tato hranice.
Stroj dává jen krok, který už v souhlasu a v péči je.
Kde není, tam nedává nic.
\end{remark}

\begin{thebibliography}{9}
\bibitem{sphere}
J.~Mikulík.
\newblock \emph{Golden Sphere}.
\newblock Přípustný krok, invariance a interface přání.
\newblock Listopad 2026.

\bibitem{principium}
J.~Mikulík.
\newblock \emph{Principium relationis}.
\newblock \texttt{notes/PrincipiumRelationis.tex}, commit \texttt{cf10e5f}.

\bibitem{lean}
J.~Mikulík.
\newblock \texttt{CausalSeed/LoveMachine.lean}.
\newblock \url{https://github.com/jendamikulik/finite-Shannon-theorem}
\end{thebibliography}

\end{document}
